\textbf{Chi tiết cài đặt}

Thí sinh cần cài đặt hai hàm sau:

\begin{lstlisting}
vector<int> solveGenius(int N, int S, const vector<int>& U, const vector<int>& V);
\end{lstlisting}

\begin{itemize}
   \item $N$: Số lượng người dân trong vương quốc ABC.
   \item $S$: Số lượng người dân tối đa có thể rời khỏi vương quốc cùng lúc trong một lượt.
   \item $U,V$: Hai mảng gồm $K=|U|=|V|$ phần tử, trong đó phần tử $U_i$ và $V_i$ ($0\le i\le K-1$) tương ứng với cặp số thứ $i$ của nhà vua.
   \item Hàm này cần trả về một mảng $A$ gồm $N$ phần tử, trong đó $A_i$ ($0\le i\le N-1$) là số nguyên dương được đánh trên chiếc áo choàng thứ $i$ mà nhà thông thái sẽ phát cho người dân thứ $i$.
   \item Hàm này được gọi \textbf{một lần duy nhất} với mỗi test.
\end{itemize}

\begin{lstlisting}
bool solveCitizen(const vector<int>& A,const vector<vector<int>>& history);
\end{lstlisting}

\begin{itemize}
   \item $A$: Một mảng gồm $N-1$ phần tử bao gồm các số được viết trên những chiếc áo choàng mà người dân tương ứng có thể nhìn thấy. Các phần tử trong dãy được sắp xếp tăng dần.
   \item $history$: Một mảng hai chiều lưu lại lịch sử của những người dân đã rời khỏi vương quốc trong các lượt trước. Trong đó:
   \begin{itemize}
      \item $|history|$ là số lượt mà nhà vua đã thực hiện trước đó.
      \item Với mọi $i$ thỏa mãn $0\le i\le |history|-1$, $history_i$ lưu trữ số được viết trên những chiếc áo choàng của những người dân đã rời khỏi vương quốc trong lượt thứ $i+1$. Các giá trị trong mảng $history_i$ được sắp xếp tăng dần.
   \end{itemize}
   \item Hàm này cần trả về \texttt{true} nếu người dân này quyết định rời khỏi vương quốc trong lượt thứ $|history|+1$ và trả về \texttt{false} trong trường hợp ngược lại.
   \item Hàm này sẽ được gọi tối đa $N^2$ lần sau khi hàm \texttt{solveGenius()} được gọi, mỗi lần gọi tương ứng với việc kiểm tra một người dân có rời khỏi vương quốc trong một lượt cụ thể không.
\end{itemize}

Trình chấm của bài toán này sẽ hoạt động như sau:
\begin{itemize}
   \item Đầu tiên, trình chấm gọi hàm \texttt{solveGenius()} trước. Thông tin của mảng $A$ từ hàm \texttt{solveGenius()} sẽ được trình chấm lưu trữ lại.
   \item Sau đó, trình chấm sẽ gọi hàm \texttt{solveCitizen()} tối đa $N^2$ lần. Cụ thể:
   \begin{itemize}
      \item Trong một lượt mô phỏng việc nhà vua đưa người dân rời khỏi vương quốc, trình chấm sẽ gọi hàm \texttt{solveCititzen()} một số lần, mỗi lần gọi là một lần kiểm tra một người dân \textbf{còn ở trong vương quốc} có muốn rời trong lượt đó không.
      \item Trình chấm sẽ thực hiện $N$ lượt tương ứng với $N$ lượt nhà vua đưa người dân rời khỏi vương quốc. Mỗi lượt, khi đã nhận về kết quả từ các lần gọi hàm \texttt{solveCititzen()} trong lượt đó, trình chấm sẽ kiểm tra tính hợp lệ của hành động này. Nếu hợp lệ, trình chấm tiếp tục xét lượt tiếp theo, hoặc dừng lại nếu đã xong $N$ lượt.
   \end{itemize}
   \item \textbf{Lưu ý: Các lần gọi hàm \texttt{solveGenius()} và \texttt{solveCitizen()} được gọi trên các chương trình hoàn toàn riêng biệt}.
   \item Cuối cùng, trình chấm tính điểm dựa trên độ tối ưu chương trình của bạn (sẽ nói rõ ở phần cách tính điểm).
\end{itemize}

\textbf{Lưu ý:}
\begin{itemize}
   \item Chương trình của bạn cần phải khai báo thư viện \texttt{"festivallib.h"} ở đầu.
   \item Trong chương trình, thí sinh được phép cài đặt các biến, các hàm toàn cục, nhưng không được phép có hàm \texttt{main()}.
   \item Chương trình của bạn không được tương tác trực tiếp với đầu vào chuẩn (stdin) và đầu ra chuẩn (stdout).
\end{itemize}

\textbf{Ràng buộc}
\begin{itemize}
   \item $2\le N\le 200$
   \item $1\le S\le N$
   \item $0\le K<\frac{N\times (N-1)}2$
   \item Đảm bảo luôn tồn tại phương án để những người dân có thể rời khỏi vương quốc
\end{itemize}